\section{Comparison with Jev on TypeSafe's published workflow evaluations}
\label{sec:jev}

TypeSafe does not publish standard benchmark results for Jev, by stated policy
\citep{typesafe_antibench}. It publishes workflow evaluations instead \citep{typesafe_evals}:
four automation tasks (Security Incidents, Agent Trace Observability, Invoice Processing and
Customer Service), each decomposed into typed questions whose answers code combines into an
action. For five example cases per workflow, the site releases the state documents, the
exact System One questions of every workflow step, the per-question answers and
probabilities of Claude Opus 5, GPT-5.6 Sol and Jev (\texttt{typesafe:v13\_snowy\_elephant}),
and reference answers from GPT-6 Astra and Claude Fable 5.1 at high thinking. TypeSafe scores
every model against the mean of the two references. These 20 cases are the only published
material on which Bobcat and Jev can be compared question by question.

\subsection{Protocol}

We fetched the four case files on 2026-09-25 (URLs and SHA-256 digests are recorded with the
results) and sent every workflow step's state and questions, unchanged, to the served Bobcat
(the merged FP8 artifact of Section~\ref{sec:engine} on one L40S): 46 requests and 418
questions, with no failure and no retry. Each request was sent once, and nothing about
Bobcat, including its prompt and temperature, was changed for or after this data. We did not
call Jev; its published answers are the only Jev output used in this work, as a comparison
row and never for training, distillation, reward or calibration.

All four models are scored by the same code. The target of a question is the modal answer of
the mean reference distribution; where the references published only values (Invoice
Processing), the two values must agree. Seventeen disputed or tied questions are excluded,
as is one question whose wording is built from each model's own earlier answers and therefore
differs between models. A Noul counts as \emph{true} above 0.5, and a Choice or Score as its
highest-probability candidate. The comparison uses the 329 questions that all four models
answered. Besides modal agreement we report the mean probability a model places on the target
and, where the references published probabilities, the mean total-variation distance to the
reference distribution. Intervals come from a bootstrap over the 20 cases, stratified by
workflow, with 10,000 resamples.

\begin{figure}[t]
\centering
\figfile{figures/vs_jev.pdf}
\caption{Question-level agreement with TypeSafe's reference consensus on the 329 questions
all four models answered in TypeSafe's 20 published workflow examples. Jev, Opus 5 and Sol
are TypeSafe's published answers; Bobcat was run once as served.}
\label{fig:vsjev}
\end{figure}

\begin{table}[t]
\centering
\small
\caption{Agreement with the reference consensus on TypeSafe's published workflow examples
(329 questions answered by all four models). ``Mass'' is the mean probability on the target
answer; TV is the mean total-variation distance to the reference distribution (lower is
better), on questions with published reference probabilities.}
\label{tab:jev}
\begin{tabular}{@{}lrrrr@{}}
\toprule
 & Bobcat & Jev & Opus 5 & Sol \\
\midrule
Security Incidents (26) & 76.9 & 80.8 & 80.8 & \textbf{88.5} \\
Agent Trace Observability (52) & 78.8 & 78.8 & \textbf{84.6} & 80.8 \\
Invoice Processing (167) & \textbf{98.8} & 97.0 & 98.2 & \textbf{98.8} \\
Customer Service (84) & \textbf{90.5} & 89.3 & 89.3 & \textbf{90.5} \\
\midrule
All questions (329) & 91.8 & 90.9 & 92.4 & \textbf{93.0} \\
Mean of the four workflows & 86.3 & 86.5 & 88.2 & \textbf{89.6} \\
Mass on the target & 0.889 & 0.850 & 0.851 & \textbf{0.914} \\
TV to the reference $\downarrow$ & 0.131 & 0.143 & \textbf{0.110} & 0.129 \\
\midrule
Noul (216) & 93.5 & 93.1 & \textbf{96.3} & 94.4 \\
Choice (87) & 90.8 & 90.8 & 90.8 & \textbf{94.3} \\
Score (26) & \textbf{80.8} & 73.1 & 65.4 & 76.9 \\
\bottomrule
\end{tabular}
\end{table}

\subsection{Results}

Bobcat and Jev agree with the reference equally often (Table~\ref{tab:jev},
Figure~\ref{fig:vsjev}): 91.8\% against 90.9\% of all questions and 86.3\% against 86.5\%
averaged over workflows, a difference of $-0.2$ points [$-3.8$, $+4.8$]. The two frontier
models score higher, but their intervals against Bobcat also include zero (Opus 5 [$-7.8$,
$+2.7$], Sol [$-10.6$, $+1.6$]). The clearest difference is in the probabilities: Bobcat places
0.889 on the target answer, Jev 0.850. Bobcat is the lowest of the four models on Security
Incidents and ties the best model on Invoice Processing and Customer Service. The 26 Score
questions are too few to interpret Bobcat's lead on that primitive.

\paragraph{Case selection.}
TypeSafe chose each workflow's five cases by how Opus 5, Sol and Jev behaved at the action
level: one case in which each of them alone differs from the other two, one in which all
three miss the reference and one in which all three agree. Each of those models therefore
carries a case chosen for its own disagreement, and Bobcat, which played no part in the
choice, does not. Without the case chosen for Jev's disagreement (256 questions), Bobcat
scores 91.0\% and Jev 90.6\% (workflow mean $-0.3$ [$-4.3$, $+5.4$]); on the two cases per
workflow that single out no model (118 questions), 89.8\% and 89.0\% ($+0.8$ [$-1.8$,
$+9.6$]). Bobcat places more probability on the target in every subset (0.889, 0.885 and
0.869 against 0.850, 0.848 and 0.839).

\paragraph{Time.}
Summing each case's step requests, sent one after another
over HTTP inside the origin host, Bobcat takes 1.89\,s per case at the median; TypeSafe
reports 0.42\,s for Jev, measured from client machines on the US West Coast with the network
included, and 20.9\,s and 24.2\,s for Opus 5 and Sol. By workflow the Bobcat-to-Jev ratio runs
from about 3 (Agent Trace Observability, 1.61 against 0.49\,s) to 19 (Invoice Processing,
8.69 against 0.45\,s). Two causes are identifiable. Our deployment runs on an L40S, about three
times slower than the H100 of Section~\ref{sec:engine} for W1. And the vLLM path recomputes
the state for every question: one Invoice Processing request carried 80 questions over a
state of about 9,000 tokens and processed 740,088 tokens in 14.2\,s, the workload for which
the state branching of Section~\ref{sec:serving} was built but which the served engine does
not yet use.

\subsection{What this comparison does not show}

The sample is small (20 cases, 329 questions) and English only, and the reference is a
consensus of two frontier models, not ground truth; TypeSafe notes that it favours OpenAI and
Anthropic models \citep{typesafe_blog}. Question-level agreement is not action accuracy.
TypeSafe also reports action accuracy on all 711 cases of the four workflows (Jev 67.8\%,
Opus 5 73.1\%, Sol 74.1\%), but neither those cases nor the code that turns answers into actions
is public, so we could not run Bobcat on that measure, and we do not place those figures next
to ours. Finally, the served artifact differs from the model behind the sealed final of
Section~\ref{sec:final} on 1.7\% of development answers. We therefore state the result as
parity with Jev on TypeSafe's published examples, not as superiority.
